\documentclass[11pt]{article}
\usepackage[a4paper,margin=27mm]{geometry}
\usepackage{amsmath,amssymb,amsthm,mathtools,microtype}
\usepackage[T1]{fontenc}
\usepackage{lmodern}
\usepackage[colorlinks=true,linkcolor=blue,urlcolor=blue,citecolor=blue]{hyperref}
\newtheorem{theorem}{Theorem}[section]
\newtheorem{proposition}[theorem]{Proposition}
\newtheorem{lemma}[theorem]{Lemma}
\newtheorem{corollary}[theorem]{Corollary}
\theoremstyle{remark}\newtheorem{remark}[theorem]{Remark}
\newcommand{\D}{\mathbb D}
\newcommand{\E}{\mathbb E}
\newcommand{\Pp}{\mathbb P}
\newcommand{\dd}{\,d}
\newcommand{\TV}{\mathrm{TV}}
\newcommand{\Rea}{\operatorname{Re}}
\newcommand{\tr}{\operatorname{tr}}
\DeclareMathOperator{\Var}{Var}
\title{Perturbing the hard-disk gas around $\beta=2$\\[3pt]\large Exact determinantal representations and uniform first responses}
\author{Research note}
\date{8 September 2026}
\begin{document}
\maketitle
\begin{abstract}
We give an exact reduced-Palm formula for temperature perturbations of the
hard-disk logarithmic gas. We prove local convergence of the first temperature
response of every fixed correlation function and hyperbolic covariance of
these first responses. An elementary total variation estimate proves the
Bergman limit when the temperature approaches $2$ at specified
$N$-dependent rates. Below $2$, the finite gas also has an exact representation
as a determinant-biased mixture of determinantal processes in a Gaussian
multiplicative chaos measure. Finally, we formulate a sufficient connected
cumulant criterion for a fixed temperature interval. That criterion is not
proved here. Neither uniqueness at any fixed $\beta\ne2$ nor the conjectured
exact intensity is claimed.
\end{abstract}

\section{Model and scope}
Put $d\mu(z)=dA(z)/\pi$ on $\D=\{|z|<1\}$ and define
\[
 S_N(\mathbf z)=\sum_{i<j}\log|z_i-z_j|,\qquad
 Z_N(b)=\int_{\D^N}e^{bS_N}\,d\mu^N,\qquad
 d\Pp_{N,b}=Z_N(b)^{-1}e^{bS_N}\,d\mu^N.
\]
Here the particles are labelled; passing to their counting measure does not
change any local assertion. Write $R_{N,b}^{(k)}$ for factorial correlation
densities with respect to $\mu^k$. Densities with respect to ordinary area are
$\rho_{N,b}^{(k)}=\pi^{-k}R_{N,b}^{(k)}$.
At $b=2$,
\[
 Z_N(2)=1,\qquad
 K_N(z,w)=\sum_{a=1}^N a(z\overline w)^{a-1},\qquad
 R_{N,2}^{(k)}(\mathbf z)=\det[K_N(z_i,z_j)]_{i,j=1}^k.
\]
Indeed, the squared norms of $1,z,\ldots,z^{N-1}$ in $L^2(\mu)$ are
$1,1/2,\ldots,1/N$, and the Vandermonde integral is $N!$ times their product.
The locally uniform kernel limit is
\[
 K(z,w)=(1-z\overline w)^{-2}.
\]
It defines the Bergman determinantal process. In this note, all temperature
derivatives are taken \emph{before} sending $N$ to infinity. No assertion
about differentiability of an unconstructed general-temperature limit is
implicit.

\section{An exact determinant-times-correction identity}
Fix $k\le N$ and distinct roots $\mathbf z=(z_1,\ldots,z_k)$ in $\D$.
Let $\Pp_{N,2}^{!\mathbf z}$ be the reduced Palm law of the $N$-particle
determinantal process at these roots; it has $N-k$ remaining particles $Y$.
Set
\[
 L_{\mathbf z}(Y)=S_{N-k}(Y)+\sum_{i=1}^k\sum_{y\in Y}\log|z_i-y|,
 \qquad \Delta_k(\mathbf z)=\prod_{i<j}(z_i-z_j).
\]
\begin{proposition}\label{palm}
For real $\varepsilon>-2$,
\begin{equation}\label{palmid}
 R_{N,2+\varepsilon}^{(k)}(\mathbf z)
 =\det[K_N(z_i,z_j)]\,|\Delta_k(\mathbf z)|^\varepsilon
 \frac{\E_{N,2}^{!\mathbf z}e^{\varepsilon L_{\mathbf z}}}
 {\E_{N,2}e^{\varepsilon S_N}}.
\end{equation}
Consequently its logarithmic correction has the Taylor germ
\begin{equation}\label{cumulants}
 \log\frac{R_{N,2+\varepsilon}^{(k)}(\mathbf z)}
 {R_{N,2}^{(k)}(\mathbf z)|\Delta_k(\mathbf z)|^\varepsilon}
 =\sum_{m\ge1}\frac{\varepsilon^m}{m!}
 \left\{\kappa_m^{!\mathbf z}(L_{\mathbf z})-\kappa_m(S_N)\right\}.
\end{equation}
The radius of this germ is not asserted to be uniform in $N$.
\end{proposition}
\begin{proof}
Integrate the density with $k$ specified coordinates. Split the pair sum
into root-root, root-remaining, and remaining-remaining terms. Dividing by
the corresponding expression at $b=2$ gives \eqref{palmid}. All integrals
are holomorphic when $\Re\varepsilon>-2$: on a compact subregion, their
absolute integrands are bounded by the same expressions with a strictly
positive Vandermonde exponent, and logarithmic derivatives are integrable.
The denominator is nonzero at zero. Taking logarithms near zero proves
\eqref{cumulants} by the definition of cumulants.
\end{proof}

The collision factor cannot be omitted from a uniform relative estimate.
Already for $N=k=2$ the relative density is a constant depending on
$\varepsilon$ times $|z_1-z_2|^\varepsilon$. It tends to zero at collision
for $\varepsilon>0$ and diverges there for $\varepsilon<0$. Thus an estimate
of the form $R_{N,2+\varepsilon}^{(2)}/R_{N,2}^{(2)}=1+O(\varepsilon)$
uniformly over all distinct points is false. One must extract the collision
factor, restrict to separated roots, or use absolute errors or local Laplace
functionals.

\section{Uniform first responses of all fixed correlations}
Write $\ell(x,y)=\log|x-y|$, and define
\[
 Q_N(x,y)=K_N(x,\mathbf z)K_N(\mathbf z,\mathbf z)^{-1}K_N(\mathbf z,y),
 \quad q_N(x)=Q_N(x,x),\quad
 L_{\mathbf z}(x)=\sum_{i=1}^k\ell(z_i,x).
\]
The reuse of $L$ for the one-particle insertion potential is restricted to
this section. The Palm kernel is $K_N-Q_N$, and $Q_N$ is an orthogonal
projection of rank $k$ on $L^2(\mu)$. In particular $\int q_N\,d\mu=k$.
Let
\[
 g_N(x)=\frac12\sum_{a=1}^N\frac{|x|^{2a}}a,
 \qquad g(x)=-\frac12\log(1-|x|^2).
\]
Use $Q,q$ for the same Palm update formed from $K$.

\begin{theorem}\label{first}
Fix a compact set of distinct root tuples in $\D^k$. Uniformly on that set,
\[
 B_N(\mathbf z):=
 \left.\partial_b\log\frac{R_{N,b}^{(k)}(\mathbf z)}
 {|\Delta_k(\mathbf z)|^b}\right|_{b=2}
 \longrightarrow B(\mathbf z),
\]
where the finite and infinite expressions are given by
\begin{align}\label{Bformula}
 B_N={}&\sum_i g_N(z_i)-\int q_N(x)g_N(x)\,d\mu(x)
          -\int L_{\mathbf z}(x)q_N(x)\,d\mu(x)\notag\\
 &+\frac12\iint\ell(x,y)
       \{q_N(x)q_N(y)-|Q_N(x,y)|^2\}\,d\mu(x)d\mu(y)\notag\\
 &+\iint\ell(x,y)\Rea\{K_N(x,y)\overline{Q_N(x,y)}\}
       \,d\mu(x)d\mu(y).
\end{align}
The expression for $B$ replaces $N$ by infinity. All displayed integrals
are absolutely convergent. In particular the first derivatives
$\partial_b R_{N,b}^{(k)}|_{b=2}$ converge locally uniformly off collisions.
\end{theorem}

\begin{lemma}\label{logkernel}
The kernels satisfy
\[
 \iint |\ell(x,y)|\,|K_N(x,y)-K(x,y)|\,d\mu(x)d\mu(y)\longrightarrow0.
\]
Also $\iint |\ell|\,|K|\,d\mu^2<\infty$.
\end{lemma}
\begin{proof}
Put $u=x\overline y$. The exact truncation identity gives
\[
 |K_N-K|\le\frac{|u|^N}{|1-u|^2}
                   +\frac{N|u|^N}{|1-u|}.
\]
On the region $\min(|x|,|y|)\le1/2$, the denominators are bounded below,
and the logarithmic diagonal singularity is uniformly integrable.
There the first term tends to zero geometrically, as does the second.
On the remaining region put $r=|x|,s=|y|$ and
$\theta=\arg x-\arg y\in[-\pi,\pi]$. Up to fixed positive constants,
\[
 |1-u|^2\ge c\{(1-rs)^2+\theta^2\},\qquad
 |\ell(x,y)|\le C(1+|\log|\theta||).
\]
The last bound holds outside the measure-zero set $\theta=0$ and follows
from $|x-y|\ge c|\theta|$ and $|x-y|\le2$. Splitting the angular integral
at $|\theta|=1-rs$ shows that its bounds for the squared and first-power
denominators are, respectively,
\[
 C\frac{1+|\log(1-rs)|}{1-rs},
 \qquad C\{1+|\log(1-rs)|^2\}.
\]
For nonnegative $f$, the radial integration identity is
\[
 \int_0^1\int_0^1 f(rs)\,r\,dr\,s\,ds
 =\int_0^1 f(v)v(-\log v)\,dv.
\]
Thus the $|1-u|^{-2}$ term without $|u|^N$ is integrable; its version with
$|u|^N$ tends to zero by dominated convergence. The remaining truncation
term is bounded by
\[
 CN\int_0^1 v^{N+1}(-\log v)
                 \{1+|\log(1-v)|^2\}\,dv
 =O\!\left(\frac{1+\log^2 N}{N}\right).
\]
For completeness, this last bound follows by splitting at $v=1/2$,
using $-\log v\le2(1-v)$ on the upper half, and then setting
$s=N(1-v)$ and using $v^N\le e^{-N(1-v)}$. The lower half decays
geometrically. These estimates prove the assertions.
\end{proof}

\begin{proof}[Proof of Theorem \ref{first}]
Radial angular averaging, or the mean-value identity for the logarithm,
gives the exact potential
\begin{equation}\label{potential}
 h_N(x):=\int\ell(x,y)K_N(y,y)\,d\mu(y)
       =-\frac{\mathcal H_N}{2}+g_N(x),
 \qquad\mathcal H_N=\sum_{a=1}^N\frac1a.
\end{equation}
Differentiate \eqref{palmid} at zero. In the Palm expectation of the pair
sum use the two-point determinant with kernel $K_N-Q_N$. Expanding that
determinant and subtracting the original pair expectation gives
\begin{align*}
 B_N={}&\sum_i h_N(z_i)-\int q_Nh_N-\int L_{\mathbf z}q_N\\
 &+\tfrac12\iint\ell(q_Nq_N-|Q_N|^2)
       +\iint\ell\Rea(K_N\overline{Q_N}).
\end{align*}
All unmarked integrals here use $\mu$. Since $\int q_N=k$, the
$-\mathcal H_N/2$ terms cancel \emph{exactly}. This is \eqref{Bformula}.

For roots in the specified compact set, the matrices
$K_N(\mathbf z,\mathbf z)$ have uniformly bounded inverses for all
sufficiently large $N$. To see this, their limits are positive definite:
the interior evaluation kernels are linearly independent, as follows by
testing against interpolation polynomials. The least eigenvalue has a
positive minimum on the compact set. Further,
$K_N(x,z_i)\to K(x,z_i)$ uniformly for $|x|\le1$ and the permitted roots.
Consequently $Q_N\to Q$ uniformly on $\overline\D^2$, with uniform bounds.
The identity $\int q=k$ follows either from projection rank or this convergence.

Now $0\le g_N\le g$ and $g\in L^1(\mu)$, so the $q_Ng_N$ integral
converges. The insertion logarithms have uniformly integrable singularities
as their roots move in a compact set, and the double logarithm lies in
$L^1(\mu^2)$. The last term converges by Lemma \ref{logkernel} and uniform
convergence of $Q_N$. All arguments are uniform in the allowed tuples.
Finally,
\[
 \left.\partial_bR_{N,b}^{(k)}\right|_2
 =\det K_N(\mathbf z,\mathbf z)
       \{\log|\Delta_k(\mathbf z)|+B_N(\mathbf z)\},
\]
which proves the final assertion.
\end{proof}

\section{Hyperbolic symmetry to first order}
Define the limiting logarithmic response
\[
 A(\mathbf z)=\log|\Delta_k(\mathbf z)|+B(\mathbf z).
\]
\begin{theorem}\label{symmetry}
For every holomorphic automorphism $\phi$ of $\D$,
\[
 A(\phi z_1,\ldots,\phi z_k)=A(z_1,\ldots,z_k).
\]
Thus the limits of all fixed correlation derivatives transform as factorial
densities of a hyperbolically invariant process. This statement concerns
first derivatives at $2$ only.
\end{theorem}
\begin{proof}
Set $a(x)=\frac12\log|\phi'(x)|$. The disk identities give
\[
 \ell(\phi x,\phi y)=\ell(x,y)+a(x)+a(y),\qquad
 g(\phi x)=g(x)-a(x).
\]
The Bergman kernel and its Palm update transform covariantly under $\phi$.
In particular $q(x)d\mu(x)$ and the two-variable measures with factors
$|Q|^2$ and $K\overline Q$ transform by substitution. The projection
identities needed below are
\[
 \int q=k,\quad \int|Q(x,y)|^2\,d\mu(y)=q(x),\quad
 \int K(x,y)\overline{Q(x,y)}\,d\mu(y)=q(x).
\]
Apply the changes of variables to \eqref{Bformula} and add the change
$(k-1)\sum_i a(z_i)$ of $\log|\Delta_k|$. The result is
\[
 A(\phi\mathbf z)-A(\mathbf z)
       =2\left\{\int q(x)a(x)\,d\mu(x)-\sum_i a(z_i)\right\}.
\]
The braces vanish by a finite-rank quadrature identity. Indeed, for a
bounded holomorphic $f$, the adjoint of multiplication by $f$ preserves
$\operatorname{span}\{K(\cdot,z_i)\}$, with eigenvalues
$\overline{f(z_i)}$. Taking its trace shows
$\int q f\,d\mu=\sum_i f(z_i)$. Take real parts and use a holomorphic
logarithm of $\phi'$; it is bounded on $\D$ for a fixed disk automorphism.
All changes of variables and traces are justified by the absolute
integrability in Theorem \ref{first} and the boundedness of $a$.
\end{proof}

\begin{corollary}\label{intensityfirst}
Locally uniformly on $\D$,
\[
 \lim_{N\to\infty}\left.\partial_b\rho_{N,b}^{(1)}(z)\right|_2
       =-\frac1{\pi(1-|z|^2)^2}.
\]
\end{corollary}
\begin{proof}
For $k=1$, Theorem \ref{symmetry} says that $A$ is constant. We compute it
at zero. There $Q_N=1$, so \eqref{Bformula} reduces to
\[
 B_N(0)=-\frac{N}{2(N+1)}+\frac12+J_N,
 \qquad J_N=\iint\ell(x,y)\Rea K_N(x,y)\,d\mu(x)d\mu(y).
\]
The constant kernel term in $J_N$ is $-1/4$. The angular expansion
\[
 \log|re^{i\theta}-s|
 =\log\max(r,s)-\sum_{n\ge1}\frac{(\min(r,s)/\max(r,s))^n}{n}\cos(n\theta)
\]
shows that the term $(n+1)(x\overline y)^n$, for $n\ge1$, contributes
$-1/[n(n+2)]$. Integration may first be done off $r=s$, followed by
integrable limiting arguments. Therefore
\[
 J_N=-\frac14-\sum_{n=1}^{N-1}\frac1{n(n+2)}
       =-1+\frac1{2N}+\frac1{2(N+1)}.
\]
It follows that $B_N(0)=-1+1/(2N)+1/(N+1)\to-1$.
Combine this with Theorem \ref{first} and the diagonal of $K$.
\end{proof}

\section{An actual determinant comparison in a shrinking window}
We use total variation with values in $[0,1]$.
\begin{theorem}\label{tv}
For $N\ge2$ and $\varepsilon\ge0$,
\[
 \|\Pp_{N,2+\varepsilon}-\Pp_{N,2}\|_{\TV}
 \le\left(\frac12+\log2\right)(N-1)\varepsilon.
\]
For $0\le\delta\le2$,
\[
 \|\Pp_{N,2-\delta}-\Pp_{N,2}\|_{\TV}
 \le1-e^{-\delta N\log N/2}\le\frac{\delta N\log N}{2}.
\]
Both bounds may of course be replaced by their minimum with $1$.
\end{theorem}
\begin{proof}
Hadamard's inequality for the Vandermonde matrix, each of whose rows has
Euclidean norm at most $\sqrt N$, yields
\[
 S_N\le M_N:=\tfrac N2\log N,\qquad T_N:=M_N-S_N\ge0.
\]
We first evaluate its mean at $2$:
\begin{equation}\label{meanenergy}
 \E_{N,2}S_N=\frac12\sum_{1\le b<a\le N}\frac1{a-b}
                  -\sum_{1\le b<a\le N}\frac1{a+b}.
\end{equation}
Here are details of the determinantal calculation. Write
$e_a(z)=\sqrt a\,z^{a-1}$, $1\le a\le N$. The angular average of
$\ell$ against independent radial densities $|e_a|^2\mu,|e_b|^2\mu$ is
$-1/[2(a+b)]$. This follows on setting $t=-\log r$ and using that these
radial variables are independent exponentials of rates $2a,2b$; the
negative logarithm of their maximum radius is their minimum.
For $a\ne b$, the angular Fourier expansion used above gives
\[
 \iint\ell(x,y)e_a(x)\overline{e_b(x)}
       \overline{e_a(y)}e_b(y)\,d\mu^2
       =-\frac{\min(a,b)}{(a+b)|a-b|}.
\]
To verify the radial factor when $a>b$, split into $r>s$ and $s>r$;
the angular coefficient is $-(s/r)^{a-b}/[2(a-b)]$ on the first region,
and its radial integral including the factor $ab$ is
$-b/[2(a+b)(a-b)]$. The other region is equal to it. For $a=b$ the
integral is $-1/(4a)$. Substitute these expressions into
$\E S_N=\frac12\iint\ell\{K_N(x,x)K_N(y,y)-|K_N(x,y)|^2\}\,d\mu^2$.
The diagonal terms cancel; the contribution of an unordered pair $a>b$
is $1/[2(a-b)]-1/(a+b)$, proving \eqref{meanenergy}.

Since $\sum_{a>b}(a-b)^{-1}=N\mathcal H_{N-1}-(N-1)$,
$\mathcal H_{N-1}\ge\log N$, and
$\sum_{b=1}^{a-1}(a+b)^{-1}<\log2$, we obtain
\[
 \E_{N,2}T_N\le\left(\tfrac12+\log2\right)(N-1).
\]
If $0<X\le1$, the probability tilted by $X/\E X$ has overlap with the
original probability at least $\E X$, because
$\min(1,X/\E X)\ge X$. Apply this with $X=e^{-\varepsilon T_N}$ to get
the first bound from $1-e^{-u}\le u$.

For the negative side, H\"older's inequality and $Z_N(0)=Z_N(2)=1$ give
$Z_N(2-\delta)\le1$. Hence
$\E_{N,2}e^{\delta T_N}=e^{\delta M_N}Z_N(2-\delta)\le e^{\delta M_N}$.
If $X\ge1$, the overlap is at least $1/\E X$. Apply this with
$X=e^{\delta T_N}$ to obtain the second bound.
\end{proof}

\begin{corollary}
The local point-process limit is Bergman if either
$\beta_N\ge2$ and $N(\beta_N-2)\to0$, or
$\beta_N\le2$ and $N\log N(2-\beta_N)\to0$.
\end{corollary}
\begin{proof}
Total variation decreases under restriction to a local window and under
every measurable map. The theorem therefore compares all bounded local
observables to those at $2$. The latter converge to their Bergman law,
for example by the local Fredholm determinant formula and convergence of
the positive projection kernels in local trace norm.
\end{proof}
These are shrinking windows; they do not prove any fixed interval of
temperatures. They also concern bounded observables, and are not asserted
to bound unbounded factorial moments without additional integrability.

\section{Below two: an exact mixture of determinantal processes}
There is a second precise interpretation of ``nearly determinantal''.
We use a standard Gaussian multiplicative chaos (GMC) construction
\cite{RV,RobertVargas}; its basic existence theorem is the only stochastic
field input in this section.

The kernel $C(z,w)=-\log|z-w|$ is positive definite as a generalized
covariance on the disk. Indeed,
\[
 C(z,w)=\underbrace{-\log\left|\frac{z-w}{1-z\overline w}\right|}_{g_D(z,w)}
       -\log|1-z\overline w|.
\]
The first term is the Dirichlet Green covariance; the second is positive
definite because it is
$\Rea\sum_{n\ge1}(z\overline w)^n/n$.
Let $X$ be a real Gaussian generalized field with covariance $C$ and, for
$0<t<2$, let
\[
 M_t(dz)={:}e^{\sqrt t X(z)}{:}\,d\mu(z)
\]
be its Wick-normalized subcritical chaos measure. This notation denotes
the limit of regularized exponentials with one half of their variance
subtracted, not an exponential of pointwise values of $X$.
The measure is finite, non-atomic, and has full support almost surely.
Its expectation is $\mu$. Off the diagonals its product moment measure is
\begin{equation}\label{gmc_moment}
 \E[M_t(dz_1)\cdots M_t(dz_N)]
    =|\Delta_N(\mathbf z)|^{-t}\,d\mu^N(\mathbf z).
\end{equation}
For example, apply the Gaussian exponential formula on pairwise separated
compact sets. Product masses are uniformly integrable in the regularization:
cross-covariances are uniformly bounded, so Gaussian comparison with
independent copies of the fields plus a common Gaussian bounds a
$(1+\eta)$th moment of their product. Choose $1+\eta<2$; the required
single-set GMC moments are uniformly bounded since $t<2$. Expectations
therefore pass to the limit. Exhaustion away from the diagonals gives the
corresponding nonnegative integral identities. No
claim that the unweighted $N$th moment of $M_t(\D)$ is finite is needed.

\begin{proposition}\label{mixture}
Define the Gram matrix and its determinant functional by
\[
 G_N(M)_{jk}=\int z^j\overline z^{,k}\,M(dz),\quad 0\le j,k<N,
 \qquad D_N(M)=N!\det G_N(M).
\]
Bias the GMC law by $D_N(M_t)/\E D_N(M_t)$ and, conditionally on $M_t$,
sample the rank-$N$ polynomial projection determinantal process in
$L^2(M_t)$. Its unconditional position law is exactly $\Pp_{N,2-t}$.
\end{proposition}
\begin{proof}
Finiteness and full support make $G_N$ finite and positive definite.
Expansion of the two Vandermonde determinants and integration gives
\[
 D_N(M)=\int|\Delta_N(\mathbf z)|^2\,M^{\otimes N}(d\mathbf z).
\]
By \eqref{gmc_moment} and non-atomicity,
\[
 \E D_N(M_t)=\int|\Delta_N|^{2-t}\,d\mu^N=Z_N(2-t)\in(0,\infty).
\]
The conditional labelled determinantal law has density
$|\Delta_N|^2/D_N(M_t)$ with respect to $M_t^{\otimes N}$. Averaging
against the biased field law cancels $D_N$ and gives exactly
$Z_N(2-t)^{-1}|\Delta_N|^{2-t}\,d\mu^N$.
\end{proof}

This is a mixture of \emph{correlation measures}: the conditional formula uses
\[
 \det[K_{N,M}(z_i,z_j)]_{i,j=1}^k\prod_{i=1}^kM(dz_i).
\]
The random measure is not an ordinary density with respect to area.
Consequently one cannot average random kernels pointwise against $dA$.
Also the mixing law depends on $N$ through a nontrivial determinant bias.
The representation by itself proves neither a local limit nor uniqueness.
For $\beta>2$ the same real Gaussian device would require a negative
covariance coefficient; it is not a positive mixture representation.

\section{A precise sufficient estimate for a fixed interval}
Let $f\in C_c^+(\D)$, $F_f(\Xi)=e^{-\langle f,\Xi\rangle}$, and define
the mixed cumulants under $\Pp_{N,2}$ by
\[
 a_{m,N}(f)=\kappa_{N,2}
   (F_f,\underbrace{S_N,\ldots,S_N}_{m\text{ copies}}),\quad m\ge1,
 \qquad a_{0,N}(f)=\E_{N,2}F_f.
\]
\begin{theorem}[Conditional criterion]\label{conditional}
Suppose that some common $0<r<2$ has the following properties for every
$f\in C_c^+(\D)$: there is $A_f<\infty$ such that
\begin{equation}\label{missing}
 \sup_N|a_{m,N}(f)|\le A_f\,m!r^{-m}\quad(m\ge1),
 \qquad a_{m,N}(f)\longrightarrow a_m(f)\quad(m\ge1).
\end{equation}
Then for every real $|\varepsilon|<r$, the Laplace functionals of
$\Pp_{N,2+\varepsilon}$ converge, and
\begin{align}\label{lapseries}
 \E_{N,2+\varepsilon}F_f
 &=\sum_{m\ge0}\frac{\varepsilon^m}{m!}a_{m,N}(f),\notag\\
 \left|\E_{N,2+\varepsilon}F_f-\E_{N,2}F_f\right|
 &\le A_f\frac{|\varepsilon|}{r-|\varepsilon|}.
\end{align}
If local point-process tightness is available at each such temperature,
the entire sequence converges there and its subsequential limit is unique.
The constants $A_f$ may depend on $f$; the radius $r$ must be common.
\end{theorem}
\begin{proof}
The Taylor coefficients of
$\E_{N,2}[F_fe^{\varepsilon S_N}]/\E_{N,2}e^{\varepsilon S_N}$ at zero
are exactly the mixed cumulants above. The bound in \eqref{missing}
makes their series converge on $|\varepsilon|<r$, uniformly in $N$ on
every smaller disk. There is no extra unproved zero-free assumption:
both the numerator and denominator are holomorphic on
$\Re\varepsilon>-2$. Multiplying the convergent series by the denominator
gives the numerator near zero, and hence throughout $|\varepsilon|<r$
by the identity theorem. On the real interval the denominator is positive,
so division proves the first assertion of \eqref{lapseries}. Summing the
geometric bound gives the second. Coefficient convergence and the common
majorant allow the $N$ limit through the series. Finally, local tightness
gives subsequential point-process limits, and equality of all their
Laplace functionals determines a unique law.
\end{proof}

At $2$ the comparison functional in \eqref{lapseries} is exactly
\[
 \E_{N,2}F_f=
 \det\!\left(I-\sqrt{1-e^{-f}}\,K_N\sqrt{1-e^{-f}}\right).
\]
Thus Theorem \ref{conditional} is a literal determinant-plus-error
criterion. For one fixed nonzero $f$, its error also proves nonemptiness
of the limit at sufficiently small $|\varepsilon|$, since the Bergman
Laplace functional is strictly less than one.
\textbf{The hypotheses \eqref{missing} are not proved in this note.}
Theorem \ref{first} verifies uniform convergence for first correlation
responses, not all orders of \eqref{missing}.

\section{Why smallness alone does not finish the argument}
There are two distinct issues. First, an $N$-uniform $O(|\varepsilon|)$
bound only places every possible limit near Bergman. It does not make
those possible limits equal. One needs convergence of the correction,
such as the convergent coefficients in Theorem \ref{conditional}.
Second, crude absolute cluster estimates lose necessary cancellations.
To see a precise integrability problem, use
\[
 d\nu(z)=\frac{d\mu(z)}{(1-|z|^2)^2},\qquad
 |\mathcal K(z,w)|=
 \frac{(1-|z|^2)(1-|w|^2)}{|1-z\overline w|^2}.
\]
At zero,
\[
 \int|\mathcal K(0,w)|\,d\nu(w)=\infty,\qquad
 \int|\mathcal K(0,w)|^2\,d\nu(w)=1.
\]
For the hyperbolic Green potential $g_D(0,w)=-\log|w|$, every $s>0$
also satisfies
\[
 \int|e^{-s g_D(0,w)}-1|\,d\nu(w)
 =\int_0^1\frac{2r(1-r^s)}{(1-r^2)^2}\,dr=\infty.
\]
Indeed the integrand is asymptotic to $s/[2(1-r)]$ at the boundary.
Making $s$ small does not remove this divergence. This rules out a direct
use of a cluster argument requiring these absolute integrals; it does not
rule out a screened or connected expansion. The exact cancellation in
\eqref{potential} is an example of the structure that such an expansion
must preserve.

The proposed intensity is, in the present normalization,
\[
 R_b^{(1)}(z)=\frac{4/b-1}{(1-|z|^2)^2}.
\]
Even proving a uniform analytic perturbation and uniqueness would not by
itself identify this exact function. Its Taylor series requires, for all
$m\ge1$,
\[
 \lim_{N\to\infty}\left.\partial_b^m R_{N,b}^{(1)}(z)\right|_2
 =\frac{(-1)^m m!\,2^{1-m}}{(1-|z|^2)^2},
\]
together with a common convergence radius, or an independent identity
that determines the density. Corollary \ref{intensityfirst} supplies the
case $m=1$. The logarithm of the conjectured density ratio is
\[
 \log\frac{4/(2+\varepsilon)-1}{1}
 =\log\frac{2-\varepsilon}{2+\varepsilon}
 =-\varepsilon-\frac{\varepsilon^3}{12}-\frac{\varepsilon^5}{80}-\cdots.
\]
Thus a further explicit test is a limiting third logarithmic response
equal to $-1/2$. It is not calculated here.

\paragraph{Relation to the literature.}
Temperature expansions from determinantal correlations have a precedent
in the one-component plasma; see \cite{Can}. That work concerns a
different background and edge regime, and is not used here to assert
a uniform remainder. The counterion disk model and its exactly soluble
coupling are discussed in \cite{Samaj}. Neither citation supplies the
missing all-order estimate for the present local limit.

\begin{thebibliography}{9}
\bibitem{RV} R. Rhodes and V. Vargas, \emph{Gaussian multiplicative chaos and
applications: a review}, Probability Surveys 11 (2014), 315--392.
\href{https://arxiv.org/abs/1305.6221}{arXiv:1305.6221}.
\bibitem{RobertVargas} R. Robert and V. Vargas, \emph{Gaussian multiplicative
chaos revisited}, Annals of Probability 38 (2010), 605--631.
\href{https://doi.org/10.1214/09-AOP490}{doi:10.1214/09-AOP490}.
\bibitem{Can} T. Can, P. J. Forrester, G. T\'ellez and P. Wiegmann,
\emph{Exact and asymptotic features of the edge density profile for the one
component plasma in two dimensions}, Journal of Statistical Physics 158
(2015), 1147--1180.
\href{https://arxiv.org/abs/1310.3130}{arXiv:1310.3130}.
\bibitem{Samaj} L. \v Samaj, \emph{Counter-Ions Near a Charged Wall:
Exact Results for Disc and Planar Geometries} (2015).
\href{https://arxiv.org/abs/1511.00882}{arXiv:1511.00882}.
\end{thebibliography}
\end{document}
